\subsection{Completed local bounds and conditional comparisons}
For each saved weighted graph, a potential is checked on every transition and
an explicit cycle attains the minimum mean. The means below use the weight
$120b_{\mathrm{new}}+L$, without division by 120. Potential ranges use the same
units. In the rational row the stored potential and weights are scaled by 19.
\begin{center}\small
\begin{tabular}{p{0.40\textwidth}r r c}
\toprule
Count-independent shares used & Mean & Range & Bound from~\eqref{eq:improved-potential}\\\midrule
Earlier catalogue, without the $012310$ floor & $480$ & $536$ & $10609/45$\\
All completed bounds in Theorem~\ref{thm:countfree-completed} & $483$ & $551$ & $157625/669$\\
Completed bounds and the proposed $012034$ floor 20 & $9259/19$ & $11301/19$ & $20965613/89059$\\
All nine target floors & $496$ & $600$ & $137941/587$\\
\bottomrule
\end{tabular}
\end{center}
The first row is retained as an earlier comparison. The second uses completed
local share theorems and gives integer bound 235, conditional on the remaining
graph interpretation and potential replay. The third adds one unproved local
share premise and still gives 235. The fourth gives 234 and retains five local
share premises. Separate integer computations verify every saved potential
inequality and every attaining cycle; these computations do not discharge the
remaining arbitrary-snake proof obligations~\cite{countfreearchive,shareclosurearchive,targetsharearchive}.

For the full target catalogue, Corollary~\ref{cor:potential234} gives
\[
4696m\le1098960+8\cdot496+600=1103528,
\qquad m\le\frac{137941}{587}<235.
\]
The strict margin at 235 is 32. A three-transition cycle with weights
$616,496,376$ attains mean 496 in this graph. Its existence does not supply
a snake construction.

The proposed $012034$ floor alone raises 1,043,097 edge weights. Its attaining
cycle has nineteen edges, and the corresponding conditional estimate is
\[
89059m\le20965613.
\]
Even an ideal potential of range zero cannot make this mean satisfy
\eqref{eq:strict234}, since
\[
111960\cdot19-227\cdot9259=25447>0.
\]
Thus completing that local theorem is useful proof progress, while additional
charge or a change to the graph or boundary argument is required to reach
234. The rational arithmetic is checked separately from the finite graph and
the pending local theorem~\cite{targetsharearchive}.
